\clearpage
\VerdierPhysicalPageAnchor{0187}%
\begingroup\small\setlength{\parskip}{0pt}\setlength{\emergencystretch}{1.5em}
\setlength{\abovedisplayskip}{4pt plus 2pt minus 2pt}%
\setlength{\belowdisplayskip}{4pt plus 2pt minus 2pt}%

\VerdierSemanticAnchor{ch3d-2-1-1-base-change}
ここで $m\colon Z\longrightarrow X$ を複体の射とする。複体 $Z'$ は、
$Z$ 上で $Z\xrightarrow{m}X$ に沿う基底変換によって得られる対象であり、次の図式の左側の縦射で与えられる：

\VerdierSemanticAnchor{ch3d-2-1-1-base-change-square}
\begin{center}
\begin{tikzcd}[column sep=4.2em,row sep=2.6em]
Z' \arrow[r] \arrow[d] & X(u,i) \arrow[d,"{p(u,i)}"] \\
Z \arrow[r,"m"'] & X
\end{tikzcd}
\end{center}
この射は、$p(w,i)\colon Z(w,i)\longrightarrow Z$ という形の複体の射と同型であり、
ここで $w$ は適当な射である。

\VerdierSemanticAnchor{ch3d-2-1-2}
\medskip\noindent\textbf{補題 2.1.2.}\quad
\textit{a) 射
$H^j(p(u,i))\colon H^j(X(u,i))\longrightarrow H^j(X)$ は
$j\neq i,i+1$ のとき同型である。}

\noindent\textit{b) 次の完全列がある：}

\VerdierSemanticAnchor{ch3d-2-1-2-exact-sequences}
\[
0\longrightarrow H^i(X(u,i))\longrightarrow H^i(X)
\longrightarrow
\frac{\operatorname{Ker}(d_X^i)+\operatorname{Im}(u)}
     {\operatorname{Im}(d_X^{i-1})+\operatorname{Im}(u)}
\longrightarrow0,
\]
\[
0\longrightarrow
\frac{\operatorname{Ker}(d_X^i)+\operatorname{Im}(u)}
     {\operatorname{Im}(d_X^{i-1})+\operatorname{Im}(u)}
\longrightarrow
\frac{X^i}{\operatorname{Im}(d_X^{i-1})+\operatorname{Im}(u)}
\longrightarrow
\frac{X^i}{\operatorname{Ker}(d_X^i)+\operatorname{Im}(u)}
\longrightarrow0,
\]
\[
0\longrightarrow
\frac{X^i}{\operatorname{Ker}(d_X^i)+\operatorname{Im}(u)}
\longrightarrow H^{i+1}(X(u,i))\longrightarrow H^{i+1}(X)
\longrightarrow0.
\]

\noindent\textit{c) $p(u,i)$ が擬同型であるための必要十分条件は}
\[
\operatorname{Im}(d_X^{i-1})+\operatorname{Im}(u)=X^i
\]
\textit{である。}
\VerdierSemanticAnchor{ch3d-2-1-2-proof}
図式追跡は読者に委ねる。

\VerdierSemanticAnchor{ch3d-2-1-3}
\medskip\noindent\textbf{定義 2.1.3.}\quad
複体の射 $m\colon Z\longrightarrow X$ が\emph{高さ $i$ の初等射}であるとは、
$\mathcal{F}\ell_1(\operatorname{comp}(\mathcal A))$ において、$p(u,i)$ 型の射と
同型であり、ここで $u$ がエピ射であることをいう。

\VerdierSemanticAnchor{ch3d-2-1-3-properties}
したがって初等射は、とくに擬同型である（(2.1.2), c)）。初等射は
$\operatorname{comp}(\mathcal A)$ のエピ射である。
$\operatorname{comp}(\mathcal A)$ のエピ射が初等射であるための必要十分条件は、
その核が高々二つの非零成分をもつ非輪状複体であることである。同じ高さの
二つの初等射の合成は
\endgroup
